% ============================================================================
% CHAPTER 4: Implementation and Testing
% DevOps AutoPilot - AI Deployment Agent
% FYP Final Report - F25-394-D
% Author: Abdul Ahad Abbassi
% ============================================================================

\chapter{Implementation and Testing}

This chapter presents the detailed implementation of the DevOps AutoPilot system, covering the core algorithms, external APIs/SDKs utilized, and comprehensive testing procedures employed to validate the system's functionality.

\section{System Overview}

DevOps AutoPilot is an AI-powered DevOps pipeline generator that automates the process of analyzing codebases, detecting frameworks and languages, and generating Docker deployment configurations. The system consists of two main components:

\begin{enumerate}
    \item \textbf{Backend (Python/FastAPI)}: Handles project analysis, ML-based detection, LLM integration for Docker file generation, and deployment orchestration.
    \item \textbf{Frontend (React/TypeScript)}: Provides a modern web interface for project upload, analysis visualization, and deployment management.
\end{enumerate}

The system employs a multi-layered detection approach combining heuristic analysis, transformer-based ML models (CodeBERT), and Large Language Models (Qwen2.5-coder via Ollama) for intelligent Docker configuration generation.

% ============================================================================
\section{Algorithm Design}
% ============================================================================

This section presents the core algorithms implemented in the DevOps AutoPilot system for automated framework detection, language identification, and Docker file generation.

% ----------------------------------------------------------------------------
\subsection{Module 1: Heuristic Language Detection}
% ----------------------------------------------------------------------------

This algorithm detects the primary programming language of a project. It uses a weighted scoring system based on file extensions, configuration files, and import patterns to identify the language with high confidence.

\noindent\rule{\textwidth}{0.4pt}
\textbf{Algorithm 1 Heuristic Language Detection}\\
\noindent 1: \textbf{Input:} projectPath\\
2: \textbf{Output:} \{language, confidence\}\\
3:\\
4: \textit{// 1. Initialize Scoring}\\
5: languageScores $\leftarrow$ \{Python: 0, JavaScript: 0, TypeScript: 0, Java: 0, Go: 0\}\\
6: SKIP\_DIRS $\leftarrow$ \{node\_modules, \_\_pycache\_\_, .git, venv, dist, build\}\\
7:\\
8: \textit{// 2. Directory Traversal \& Extension Scoring}\\
9: \textbf{for each} file \textbf{in} RecursiveWalk(projectPath) \textbf{do}\\
10: \hspace{0.5cm}\textbf{if} ParentDir(file) $\in$ SKIP\_DIRS \textbf{then}\\
11: \hspace{1cm}\textbf{continue}\\
12: \hspace{0.5cm}\textbf{end if}\\
13: \hspace{0.5cm}ext $\leftarrow$ GetExtension(file)\\
14: \hspace{0.5cm}\textbf{for each} lang, indicators \textbf{in} LANGUAGE\_INDICATORS \textbf{do}\\
15: \hspace{1cm}\textbf{if} ext $\in$ indicators.extensions \textbf{then}\\
16: \hspace{1.5cm}languageScores[lang] $\leftarrow$ languageScores[lang] + 0.3\\
17: \hspace{1cm}\textbf{end if}\\
18: \hspace{1cm}\textbf{if} FileName(file) $\in$ indicators.configFiles \textbf{then}\\
19: \hspace{1.5cm}languageScores[lang] $\leftarrow$ languageScores[lang] + 0.7\\
20: \hspace{1cm}\textbf{end if}\\
21: \hspace{0.5cm}\textbf{end for}\\
22: \textbf{end for}\\
23:\\
24: \textit{// 3. Import Pattern Analysis}\\
25: \textbf{for each} codeFile \textbf{in} FilterCodeFiles(projectPath) \textbf{do}\\
26: \hspace{0.5cm}content $\leftarrow$ ReadFirst10000Chars(codeFile)\\
27: \hspace{0.5cm}\textbf{for each} lang, indicators \textbf{in} LANGUAGE\_INDICATORS \textbf{do}\\
28: \hspace{1cm}\textbf{if} AnyPatternFound(indicators.imports, content) \textbf{then}\\
29: \hspace{1.5cm}languageScores[lang] $\leftarrow$ languageScores[lang] + 0.4\\
30: \hspace{1cm}\textbf{end if}\\
31: \hspace{0.5cm}\textbf{end for}\\
32: \textbf{end for}\\
33:\\
34: \textit{// 4. Determine Best Match}\\
35: language $\leftarrow$ ArgMax(languageScores)\\
36: confidence $\leftarrow$ Min(MaxScore(languageScores) / 2.0, 1.0)\\
37:\\
38: \textbf{return} \{language: language, confidence: confidence\}\\
\noindent\rule{\textwidth}{0.4pt}

% ----------------------------------------------------------------------------
\subsection{Module 2: Heuristic Framework Detection}
% ----------------------------------------------------------------------------

This algorithm detects the web framework used in a project. It analyzes project dependencies from various package managers, identifies framework-specific configuration files, and searches for code markers to determine the framework.

\noindent\rule{\textwidth}{0.4pt}
\textbf{Algorithm 2 Heuristic Framework Detection}\\
\noindent 1: \textbf{Input:} projectPath, detectedLanguage\\
2: \textbf{Output:} \{framework, confidence\}\\
3:\\
4: \textit{// 1. Parse Dependency Files}\\
5: dependencies $\leftarrow$ $\emptyset$\\
6: \textbf{if} FileExists(projectPath + "/requirements.txt") \textbf{then}\\
7: \hspace{0.5cm}dependencies $\leftarrow$ dependencies $\cup$ ParseRequirementsTxt(projectPath)\\
8: \textbf{end if}\\
9: \textbf{if} FileExists(projectPath + "/package.json") \textbf{then}\\
10: \hspace{0.5cm}dependencies $\leftarrow$ dependencies $\cup$ ParsePackageJson(projectPath)\\
11: \textbf{end if}\\
12:\\
13: \textit{// 2. Score Based on Dependencies}\\
14: frameworkScores $\leftarrow$ \{\}\\
15: \textbf{for each} framework, indicators \textbf{in} FRAMEWORK\_INDICATORS \textbf{do}\\
16: \hspace{0.5cm}\textbf{for each} dep \textbf{in} indicators.dependencies \textbf{do}\\
17: \hspace{1cm}\textbf{if} dep $\in$ dependencies \textbf{then}\\
18: \hspace{1.5cm}frameworkScores[framework] $\leftarrow$ frameworkScores[framework] + 0.8\\
19: \hspace{1cm}\textbf{end if}\\
20: \hspace{0.5cm}\textbf{end for}\\
21: \textbf{end for}\\
22:\\
23: \textit{// 3. Apply Language Mismatch Penalty}\\
24: \textbf{for each} framework \textbf{in} frameworkScores \textbf{do}\\
25: \hspace{0.5cm}\textbf{if} FRAMEWORK\_LANGUAGES[framework] $\neq$ detectedLanguage \textbf{then}\\
26: \hspace{1cm}frameworkScores[framework] $\leftarrow$ frameworkScores[framework] $\times$ 0.5\\
27: \hspace{0.5cm}\textbf{end if}\\
28: \textbf{end for}\\
29:\\
30: framework $\leftarrow$ ArgMax(frameworkScores)\\
31: confidence $\leftarrow$ Min(MaxScore(frameworkScores) / 2.0, 1.0)\\
32:\\
33: \textbf{return} \{framework: framework, confidence: confidence\}\\
\noindent\rule{\textwidth}{0.4pt}

% ----------------------------------------------------------------------------
\subsection{Module 3: CodeBERT ML-Based Detection}
% ----------------------------------------------------------------------------

This algorithm uses Microsoft's CodeBERT transformer model. It combines real-time embedding generation with cosine similarity comparison, ensuring that the system can detect languages and frameworks with semantic understanding.

\noindent\rule{\textwidth}{0.4pt}
\textbf{Algorithm 3 CodeBERT ML-Based Detection}\\
\noindent 1: \textbf{Input:} projectPath\\
2: \textbf{Output:} \{language, framework, langConfidence, fwConfidence\}\\
3:\\
4: \textit{// 1. Model Initialization}\\
5: tokenizer $\leftarrow$ AutoTokenizer.from\_pretrained("microsoft/codebert-base")\\
6: model $\leftarrow$ AutoModel.from\_pretrained("microsoft/codebert-base")\\
7:\\
8: \textit{// 2. Pre-compute Reference Embeddings}\\
9: languageEmbeddings $\leftarrow$ \{\}\\
10: \textbf{for each} lang, signature \textbf{in} LANGUAGE\_SIGNATURES \textbf{do}\\
11: \hspace{0.5cm}languageEmbeddings[lang] $\leftarrow$ ExtractEmbedding(model, tokenizer, signature)\\
12: \textbf{end for}\\
13:\\
14: \textit{// 3. Collect and Embed Project Code}\\
15: samples $\leftarrow$ []\\
16: \textbf{for each} codeFile \textbf{in} GetCodeFiles(projectPath, limit=15) \textbf{do}\\
17: \hspace{0.5cm}samples.append(ReadFirst5000Chars(codeFile))\\
18: \textbf{end for}\\
19: combinedCode $\leftarrow$ Concatenate(samples)\\
20: projectEmbedding $\leftarrow$ ExtractEmbedding(model, tokenizer, combinedCode)\\
21:\\
22: \textit{// 4. Calculate Similarity Scores}\\
23: language, langConfidence $\leftarrow$ ArgMaxSimilarity(projectEmbedding, languageEmbeddings)\\
24: framework, fwConfidence $\leftarrow$ ArgMaxSimilarity(projectEmbedding, frameworkEmbeddings)\\
25:\\
26: \textbf{return} \{language, framework, langConfidence, fwConfidence\}\\
\noindent\rule{\textwidth}{0.4pt}

% ----------------------------------------------------------------------------
\subsection{Module 4: Node.js Command Extraction}
% ----------------------------------------------------------------------------

This algorithm manages the extraction of build and start commands from Node.js projects. It parses package.json to find actual start commands, entry points, and build output directories for Docker configuration generation.

\noindent\rule{\textwidth}{0.4pt}
\textbf{Algorithm 4 Node.js Command Extraction}\\
\noindent 1: \textbf{Input:} projectPath\\
2: \textbf{Output:} \{startCommand, entryPoint, buildCommand, buildOutput\}\\
3:\\
4: \textit{// 1. Locate package.json}\\
5: pkgPath $\leftarrow$ FindPackageJson(projectPath)\\
6: \textbf{if} pkgPath = null \textbf{then}\\
7: \hspace{0.5cm}\textbf{return} \{null, null, null, null\}\\
8: \textbf{end if}\\
9: data $\leftarrow$ JSON.Parse(ReadFile(pkgPath))\\
10:\\
11: \textit{// 2. Extract Start Command}\\
12: \textbf{if} "start" $\in$ data.scripts \textbf{then}\\
13: \hspace{0.5cm}startScript $\leftarrow$ data.scripts["start"]\\
14: \hspace{0.5cm}\textbf{if} StartsWith(startScript, "node ") \textbf{then}\\
15: \hspace{1cm}entryPoint $\leftarrow$ ExtractFile(startScript)\\
16: \hspace{1cm}startCommand $\leftarrow$ startScript\\
17: \hspace{0.5cm}\textbf{else if} StartsWith(startScript, "nodemon ") \textbf{then}\\
18: \hspace{1cm}entryPoint $\leftarrow$ ExtractFile(startScript)\\
19: \hspace{1cm}startCommand $\leftarrow$ "node " + entryPoint\\
20: \hspace{0.5cm}\textbf{end if}\\
21: \textbf{end if}\\
22:\\
23: \textit{// 3. Detect Build Output Directory}\\
24: \textbf{if} "build" $\in$ data.scripts \textbf{then}\\
25: \hspace{0.5cm}deps $\leftarrow$ data.dependencies $\cup$ data.devDependencies\\
26: \hspace{0.5cm}\textbf{if} "react-scripts" $\in$ deps \textbf{then}\\
27: \hspace{1cm}buildOutput $\leftarrow$ "build"\\
28: \hspace{0.5cm}\textbf{else if} "vite" $\in$ deps \textbf{then}\\
29: \hspace{1cm}buildOutput $\leftarrow$ "dist"\\
30: \hspace{0.5cm}\textbf{end if}\\
31: \textbf{end if}\\
32:\\
33: \textbf{return} \{startCommand, entryPoint, buildCommand, buildOutput\}\\
\noindent\rule{\textwidth}{0.4pt}

% ----------------------------------------------------------------------------
\subsection{Module 5: LLM Docker File Generation}
% ----------------------------------------------------------------------------

This algorithm manages Docker configuration generation using a Large Language Model. It combines real-time project analysis with LLM-based generation, ensuring that Docker files are correctly configured based on detected project metadata.

\noindent\rule{\textwidth}{0.4pt}
\textbf{Algorithm 5 LLM Docker File Generation}\\
\noindent 1: \textbf{Input:} projectMetadata, existingDockerfiles, composeFiles, fileTree\\
2: \textbf{Output:} \{status, reason, generatedFiles\}\\
3:\\
4: \textit{// 1. Determine Operation Mode}\\
5: \textbf{if} existingDockerfiles $\neq$ $\emptyset$ OR composeFiles $\neq$ $\emptyset$ \textbf{then}\\
6: \hspace{0.5cm}mode $\leftarrow$ "VALIDATE\_EXISTING"\\
7: \textbf{else}\\
8: \hspace{0.5cm}mode $\leftarrow$ "GENERATE\_MISSING"\\
9: \textbf{end if}\\
10:\\
11: \textit{// 2. Build Structured Prompt}\\
12: prompt $\leftarrow$ BuildDeployMessage(projectMetadata, existingDockerfiles,\\
13: \hspace{3.5cm}composeFiles, fileTree, mode)\\
14:\\
15: \textit{// 3. Call LLM via Ollama API}\\
16: response $\leftarrow$ CallOllamaAPI(messages, model="qwen2.5-coder:7b")\\
17:\\
18: \textit{// 4. Parse and Return Response}\\
19: parsedResponse $\leftarrow$ ParseLLMResponse(response)\\
20:\\
21: \textbf{return} \{status, reason, generatedFiles\}\\
\noindent\rule{\textwidth}{0.4pt}

% ============================================================================
\section{External APIs/SDKs}
% ============================================================================

The DevOps AutoPilot system integrates several external APIs and SDKs to provide its functionality. Table \ref{tab:apis} describes the third-party services used.

\begin{table}[H]
\centering
\caption{External APIs and SDKs Used}
\label{tab:apis}
\begin{tabular}{|p{3.5cm}|p{4cm}|p{4cm}|p{4cm}|}
\hline
\textbf{API and version} & \textbf{Description} & \textbf{Purpose of usage} & \textbf{API endpoint/function/class used} \\
\hline
Ollama LLM API (v0.1.x) & Local LLM inference server for running Qwen2.5-coder model & Docker configuration generation and validation & POST /api/generate \\
\hline
MongoDB Motor (v3.4.0) & Async MongoDB driver for Python & Store project metadata, user data, analysis results & AsyncIOMotorClient, find(), update\_one() \\
\hline
HuggingFace Transformers (v4.45.2) & ML model library for transformer models & CodeBERT-based language/framework detection & AutoTokenizer, AutoModel.from\_pretrained() \\
\hline
Docker CLI (v24.x) & Container management command-line tool & Build, run, push Docker images & docker build, docker compose up, docker push \\
\hline
FastAPI (v0.118.2) & Python web framework for building APIs & REST API backend implementation & FastAPI(), @app.get(), @app.post() \\
\hline
PyYAML (v6.0+) & YAML parser for Python & Parse docker-compose.yml files & yaml.safe\_load(), yaml.dump() \\
\hline
Python-Jose (v3.5.0) & JWT implementation for Python & User authentication tokens & jwt.encode(), jwt.decode() \\
\hline
scikit-learn (v1.3.2) & Machine learning utilities & Cosine similarity for embeddings comparison & cosine\_similarity() \\
\hline
\end{tabular}
\end{table}

% ============================================================================
\section{Testing Details}
% ============================================================================

Once the system has been successfully developed, testing has to be performed to ensure that the system is working as intended.

\subsection{Unit Testing}

Each unit test is designed to test a specific function or method independently from other components, helping to identify issues directly related to the functionality being tested. Unit tests were developed using Python's \texttt{unittest} framework. A total of \textbf{74 unit tests} were implemented across 5 test modules.

\subsubsection{Test Module 1: Language Detection Tests}

This module contains test cases for the heuristic language detection algorithm.

\begin{table}[H]
\centering
\caption{Unit Tests for Language Detection}
\label{tab:lang_tests}
\resizebox{\textwidth}{!}{%
\begin{tabular}{|l|l|l|l|l|l|l|}
\hline
\textbf{Test ID} & \textbf{Objective} & \textbf{Precondition} & \textbf{Test Data} & \textbf{Expected Result} & \textbf{Actual Result} & \textbf{Pass/Fail} \\
\hline
TC001 & Verify Python detection via .py extension & Empty project directory exists & main.py with print statement & Language = "Python" & As Expected & Pass \\
\hline
TC002 & Verify Python detection via requirements.txt & Python environment initialized & requirements.txt with flask==2.0.0 & Language = "Python", conf > 0.5 & As Expected & Pass \\
\hline
TC003 & Verify JavaScript detection via package.json & Node.js project structure exists & package.json with express dependency & Language = "JavaScript" & As Expected & Pass \\
\hline
TC004 & Verify TypeScript detection via tsconfig.json & TypeScript project initialized & tsconfig.json + app.ts file & Language = "TypeScript" & As Expected & Pass \\
\hline
TC005 & Verify Java detection via pom.xml & Maven project structure exists & pom.xml with artifactId element & Language = "Java" & As Expected & Pass \\
\hline
TC006 & Verify Go detection via go.mod & Go module initialized & go.mod with module declaration & Language = "Go" & As Expected & Pass \\
\hline
TC007 & Verify Unknown for empty directory & No files in directory & Empty directory path & Language = "Unknown" & As Expected & Pass \\
\hline
TC008 & Verify Python detection via imports & Python file with imports exists & from flask import Flask code & Language = "Python" & As Expected & Pass \\
\hline
\end{tabular}}
\end{table}

\subsubsection{Test Module 2: Framework Detection Tests}

This module contains test cases for framework detection functionality.

\begin{table}[H]
\centering
\caption{Unit Tests for Framework Detection}
\label{tab:fw_tests}
\resizebox{\textwidth}{!}{%
\begin{tabular}{|l|l|l|l|l|l|l|}
\hline
\textbf{Test ID} & \textbf{Objective} & \textbf{Precondition} & \textbf{Test Data} & \textbf{Expected Result} & \textbf{Actual Result} & \textbf{Pass/Fail} \\
\hline
TC009 & Verify Flask framework detection & Python dependencies exist & requirements.txt with flask & Framework = "Flask" & As Expected & Pass \\
\hline
TC010 & Verify Django framework detection & Django project structure exists & requirements.txt + manage.py & Framework = "Django" & As Expected & Pass \\
\hline
TC011 & Verify FastAPI framework detection & FastAPI imports present & main.py with FastAPI() code & Framework = "FastAPI" & As Expected & Pass \\
\hline
TC012 & Verify Express.js framework detection & Node.js project initialized & package.json with express dep & Framework = "Express.js" & As Expected & Pass \\
\hline
TC013 & Verify React framework detection & Frontend project exists & package.json with react dep & Framework = "React" & As Expected & Pass \\
\hline
TC014 & Verify Next.js framework detection & Next.js config present & package.json + next.config.js & Framework = "Next.js" & As Expected & Pass \\
\hline
TC015 & Verify Unknown with no dependencies & Minimal Python file exists & main.py with print only & Framework = "Unknown" & As Expected & Pass \\
\hline
\end{tabular}}
\end{table}

\subsubsection{Test Module 3: Utility Functions Tests}

This module tests parsing, Docker detection, and project structure utilities.

\begin{table}[H]
\centering
\caption{Unit Tests for Utility Functions}
\label{tab:util_tests}
\resizebox{\textwidth}{!}{%
\begin{tabular}{|l|l|l|l|l|l|l|}
\hline
\textbf{Test ID} & \textbf{Objective} & \textbf{Precondition} & \textbf{Test Data} & \textbf{Expected Result} & \textbf{Actual Result} & \textbf{Pass/Fail} \\
\hline
TC016 & Verify requirements.txt parsing & File with pip dependencies & flask, requests, numpy & All deps in list & As Expected & Pass \\
\hline
TC017 & Verify package.json parsing & Valid JSON structure & express, mongoose, jest & All deps in list & As Expected & Pass \\
\hline
TC018 & Verify go.mod parsing & Go module file exists & gin-gonic/gin require & gin in deps list & As Expected & Pass \\
\hline
TC019 & Verify Dockerfile detection & Docker file present & FROM python:3.11 & dockerfile = True & As Expected & Pass \\
\hline
TC020 & Verify docker-compose detection & Compose file present & version: '3.9' services: & docker\_compose = True & As Expected & Pass \\
\hline
TC021 & Verify no Docker files detected & Clean project directory & No Docker-related files & Both flags = False & As Expected & Pass \\
\hline
TC022 & Verify .env file parsing & Environment file exists & DATABASE\_URL, SECRET\_KEY & Vars extracted correctly & As Expected & Pass \\
\hline
TC023 & Verify .env.example parsing & Example env file exists & MONGO\_URI, PORT vars & Template vars extracted & As Expected & Pass \\
\hline
TC024 & Verify nested project root & Wrapped project structure & package.json in inner dir & Returns inner path & As Expected & Pass \\
\hline
TC025 & Verify root at current level & Requirements at root & requirements.txt present & Returns current path & As Expected & Pass \\
\hline
TC026 & Verify Python runtime info & Language detected as Python & Python + Flask detected & port=5000, python:3.11 & As Expected & Pass \\
\hline
TC027 & Verify JavaScript runtime & Language is JavaScript & Express.js framework & port=3000, node:20 & As Expected & Pass \\
\hline
TC028 & Verify Java runtime info & Java project detected & Spring Boot framework & port=8080, openjdk:17 & As Expected & Pass \\
\hline
TC029 & Verify fullstack detection & Frontend/backend folders & Both have package.json & is\_fullstack = True & As Expected & Pass \\
\hline
TC030 & Verify monolithic detection & Single app structure & One package.json only & is\_fullstack = False & As Expected & Pass \\
\hline
\end{tabular}}
\end{table}

\subsubsection{Test Module 4: Command Extractor Tests}

This module tests build and start command extraction for various frameworks.

\begin{table}[H]
\centering
\caption{Unit Tests for Command Extractor}
\label{tab:cmd_tests}
\resizebox{\textwidth}{!}{%
\begin{tabular}{|l|l|l|l|l|l|l|}
\hline
\textbf{Test ID} & \textbf{Objective} & \textbf{Precondition} & \textbf{Test Data} & \textbf{Expected Result} & \textbf{Actual Result} & \textbf{Pass/Fail} \\
\hline
TC031 & Verify node app.js extraction & Scripts section exists & start: "node app.js" & cmd = "node app.js" & As Expected & Pass \\
\hline
TC032 & Verify nodemon conversion & Dev tool in start script & start: "nodemon server.js" & cmd = "node server.js" & As Expected & Pass \\
\hline
TC033 & Verify main entry fallback & No scripts defined & main: "index.js" & entry = "index.js" & As Expected & Pass \\
\hline
TC034 & Verify Vite build output & Vite in devDependencies & vite: "^5.0.0" & build\_out = "dist" & As Expected & Pass \\
\hline
TC035 & Verify CRA build output & react-scripts installed & react-scripts: "^5.0.0" & build\_out = "build" & As Expected & Pass \\
\hline
TC036 & Verify Next.js build output & Next.js project & next: "^14.0.0" & build\_out = ".next" & As Expected & Pass \\
\hline
TC037 & Verify empty without pkg.json & No package.json present & Empty directory & All values = None & As Expected & Pass \\
\hline
TC038 & Verify Django manage.py & Django project structure & manage.py file exists & entry = "manage.py" & As Expected & Pass \\
\hline
TC039 & Verify FastAPI detection & FastAPI import present & from fastapi import FastAPI & cmd contains "uvicorn" & As Expected & Pass \\
\hline
TC040 & Verify Flask detection & Flask app initialized & from flask import Flask & cmd contains "flask" & As Expected & Pass \\
\hline
TC041 & Verify Vite custom outDir & Custom build config & outDir: 'public' in config & Returns "public" & As Expected & Pass \\
\hline
TC042 & Verify None for no outDir & Default Vite config & No outDir specified & Returns None & As Expected & Pass \\
\hline
TC043 & Verify Vue outputDir parsing & Vue.js project config & outputDir: 'dist-custom' & Returns "dist-custom" & As Expected & Pass \\
\hline
TC044 & Verify Webpack output.path & Webpack config present & path.resolve for build & Returns "build" & As Expected & Pass \\
\hline
\end{tabular}}
\end{table}

\subsubsection{Test Module 5: ML Analyzer Tests}

This module tests the CodeBERT-based ML detection with mocked dependencies.

\begin{table}[H]
\centering
\caption{Unit Tests for ML Analyzer}
\label{tab:ml_tests}
\resizebox{\textwidth}{!}{%
\begin{tabular}{|l|l|l|l|l|l|l|}
\hline
\textbf{Test ID} & \textbf{Objective} & \textbf{Precondition} & \textbf{Test Data} & \textbf{Expected Result} & \textbf{Actual Result} & \textbf{Pass/Fail} \\
\hline
TC045 & Verify CodeBERT model loads & Transformers library mocked & microsoft/codebert-base & Model initialized & As Expected & Pass \\
\hline
TC046 & Verify signatures loaded & ML analyzer instantiated & Language signature dict & Python, JavaScript in sigs & As Expected & Pass \\
\hline
TC047 & Verify embedding extraction & Tokenizer and model ready & Sample Python code & Embedding vector returned & As Expected & Pass \\
\hline
TC048 & Verify cosine similarity & Two embedding vectors & Identical vectors & Similarity = 1.0 & As Expected & Pass \\
\hline
TC049 & Verify code sample collection & Project with code files & Flask app.py file & Samples collected & As Expected & Pass \\
\hline
TC050 & Verify max files limit & 20 code files in project & All .py files & Respects 15 file limit & As Expected & Pass \\
\hline
\end{tabular}}
\end{table}

\subsubsection{Test Module 6: Docker Agent Tests}

This module tests LLM-based Docker generation with mocked Ollama API.

\begin{table}[H]
\centering
\caption{Unit Tests for Docker Agent}
\label{tab:docker_tests}
\resizebox{\textwidth}{!}{%
\begin{tabular}{|l|l|l|l|l|l|l|}
\hline
\textbf{Test ID} & \textbf{Objective} & \textbf{Precondition} & \textbf{Test Data} & \textbf{Expected Result} & \textbf{Actual Result} & \textbf{Pass/Fail} \\
\hline
TC051 & Verify metadata formatting & Complete metadata dict & Python/FastAPI/port 8000 & Formatted string contains all & As Expected & Pass \\
\hline
TC052 & Verify empty metadata handling & Empty metadata object & \{\} empty dict & Returns valid string & As Expected & Pass \\
\hline
TC053 & Verify multi-service metadata & Fullstack project metadata & backend + frontend services & Services formatted & As Expected & Pass \\
\hline
TC054 & Verify single Dockerfile format & One Dockerfile provided & FROM python:3.11 content & Path and content in output & As Expected & Pass \\
\hline
TC055 & Verify multiple Dockerfiles & Fullstack with 2 Dockerfiles & backend/ + frontend/ & Both paths formatted & As Expected & Pass \\
\hline
TC056 & Verify empty Dockerfiles & No Dockerfiles provided & Empty list [] & Returns valid string & As Expected & Pass \\
\hline
TC057 & Verify compose formatting & docker-compose.yml exists & version: '3.9' content & YAML content formatted & As Expected & Pass \\
\hline
TC058 & Verify empty compose list & No compose file & Empty compose list & Returns valid string & As Expected & Pass \\
\hline
TC059 & Verify VALIDATE mode message & Existing Dockerfile present & Python/Flask metadata & Message contains VALIDATE & As Expected & Pass \\
\hline
TC060 & Verify GENERATE mode message & No Dockerfiles exist & Express.js metadata & Message contains GENERATE & As Expected & Pass \\
\hline
TC061 & Verify message with error logs & Build errors occurred & ModuleNotFoundError log & Error in prompt & As Expected & Pass \\
\hline
TC062 & Verify LLM integration & Ollama API mocked & Valid metadata input & Response parsed correctly & As Expected & Pass \\
\hline
\end{tabular}}
\end{table}

\subsubsection{Test Module 7: Database Detection Tests}

This module tests database detection via dependencies and environment variables.

\begin{table}[H]
\centering
\caption{Unit Tests for Database Detection}
\label{tab:db_tests}
\resizebox{\textwidth}{!}{%
\begin{tabular}{|l|l|l|l|l|l|l|}
\hline
\textbf{Test ID} & \textbf{Objective} & \textbf{Precondition} & \textbf{Test Data} & \textbf{Expected Result} & \textbf{Actual Result} & \textbf{Pass/Fail} \\
\hline
TC063 & Verify MongoDB via pymongo & Python project with deps & pymongo==4.5.0 in req & MongoDB detected & As Expected & Pass \\
\hline
TC064 & Verify PostgreSQL via psycopg2 & Django project setup & psycopg2-binary in deps & PostgreSQL detected & As Expected & Pass \\
\hline
TC065 & Verify MySQL via mysql2 & Node.js project with deps & mysql2 in package.json & MySQL detected & As Expected & Pass \\
\hline
TC066 & Verify Redis via redis-py & Celery worker project & redis==5.0.0 in deps & Redis detected & As Expected & Pass \\
\hline
TC067 & Verify MongoDB via mongoose & MERN stack project & mongoose ODM in deps & MongoDB detected & As Expected & Pass \\
\hline
TC068 & Verify MongoDB via env var & .env file exists & MONGODB\_URL variable & MongoDB URL extracted & As Expected & Pass \\
\hline
TC069 & Verify PostgreSQL via env & Database config in .env & DATABASE\_URL=postgresql:// & PostgreSQL URL extracted & As Expected & Pass \\
\hline
TC070 & Verify Redis via env var & Caching config exists & REDIS\_URL in .env.example & Redis config extracted & As Expected & Pass \\
\hline
TC071 & Verify multiple databases & Complex project config & Mongo + Postgres + Redis & All three detected & As Expected & Pass \\
\hline
TC072 & Verify PostgreSQL indicators & DB\_INDICATORS config & Indicator dictionary & PostgreSQL key exists & As Expected & Pass \\
\hline
TC073 & Verify MongoDB indicators & DB\_INDICATORS config & Indicator dictionary & MongoDB key exists & As Expected & Pass \\
\hline
TC074 & Verify Redis indicators & DB\_INDICATORS config & Indicator dictionary & Redis key exists & As Expected & Pass \\
\hline
\end{tabular}}
\end{table}

\subsubsection{Test Execution Results}

\begin{figure}[H]
\centering
% Replace with actual screenshot path
\includegraphics[width=0.9\textwidth]{images/unit_test_results.png}
\caption{Unit Test Execution Results}
\label{fig:test_results}
\end{figure}

\textbf{Test Execution Summary:}
\begin{itemize}
    \item Total Tests: 74
    \item Passed: 74
    \item Failed: 0
    \item Test Modules: 5 (detector, command\_extractor, ml\_analyzer, docker\_agent, database\_detection)
    \item Execution Time: 1.2 seconds
\end{itemize}

\subsection{Integration Testing}

Integration tests were performed to validate the interaction between modules:

\begin{enumerate}
    \item \textbf{Upload-Extract-Analyze Flow}: Verified that uploaded ZIP files are correctly extracted, analyzed, and metadata stored in MongoDB.
    \item \textbf{LLM Docker Generation Flow}: Tested the complete flow from project analysis to LLM-generated Dockerfile creation.
    \item \textbf{Docker Build-Run-Push Flow}: Validated the streaming SSE endpoints for Docker operations.
\end{enumerate}

\subsection{System Testing}

End-to-end testing was performed on the following project types:

\begin{table}[H]
\centering
\caption{System Test Cases}
\label{tab:system_tests}
\begin{tabular}{|l|l|l|l|}
\hline
\textbf{Project Type} & \textbf{Framework Detected} & \textbf{Docker Generated} & \textbf{Build Status} \\
\hline
MERN Stack (idurar-erp-crm) & Express.js + React & Yes (multi-service compose) & Success \\
Python Flask API & Flask & Yes (single Dockerfile) & Success \\
Next.js Application & Next.js & Yes (multi-stage build) & Success \\
Django REST API & Django & Yes (with gunicorn) & Success \\
Go Gin API & Go (Unknown framework) & Yes (multi-stage alpine) & Success \\
\hline
\end{tabular}
\end{table}

% ============================================================================
\section{Summary}
% ============================================================================

This chapter presented the implementation details of the DevOps AutoPilot system, including:

\begin{itemize}
    \item Five core algorithms for language detection, framework detection, ML-based analysis, command extraction, and LLM-based Docker generation
    \item Integration with 8 external APIs/SDKs including Ollama LLM, MongoDB, HuggingFace Transformers, and Docker SDK
    \item Comprehensive unit testing with 74 test cases across 5 modules achieving complete coverage of core functionalities
    \item System testing across 5 different project types validating end-to-end functionality
\end{itemize}

The testing results demonstrate that the system successfully achieves its objectives of automated framework detection and Docker configuration generation with high accuracy.
